\chapter*{后记}\addcontentsline{toc}{chapter}{后记}\markboth{后记}{后记}
% 原书第467–468页

2024年8月，这本书大量引用的《北京纪事》第3卷由荷兰的博睿（Brill）出版社出版。当时一个中西交往史研究的e-mail群还通知大家，《北京纪事》的主要翻译者和主编鲁保禄（Paul Anthony Rule）将在12月举办一个线上活动，介绍新卷中的内容。Paul是中西传教史研究领域内的元老，二十年前我刚到美国读书时就读过他的论著。直到我在写博士论文，大量用到《北京纪事》的第1、2卷时，才发现Paul还活跃在研究的第一线。在2019年5月完成博士论文后，我在网上找到了Paul的电邮地址，联系了他，介绍了我博士论文中的观点，并得到了他的回复指点。后来，他还把《北京纪事》第3卷的word 草稿发给了我。他告诉我第3卷出版尚待时日。如果我发现有需要引用的内容，可以按照原始未刊档案方式引用。2020年春，基于我博士论文写出的英文书通过了各种外审，我在拿到了出版合同以后询问Paul是否愿意看看底稿。Paul同意了。一个月后，他给我发来了电邮回复。出乎意料的是，放在附件中的评论长达1万字。而且在电邮的正文中，他还给我写了一小段让我印象深刻的话：


% 原书第468页

\begin{bookquote}I found your manuscript very thought provoking. I have given much time to it and the result is the attached "Comments". I hope you do not find them too confronting. I am accustomed with my friends to express myself somewhat more frankly than perhaps Chinese politeness would allow, but I do so because I think your material is so interesting and wide ranging. And I hope you regard them as friendly. Don't hesitate to reply as frankly.\end{bookquote}

% 原书第468页

\begin{bookquote}你的稿子让我想到了很多东西。我花了许多时间在上面，想法都在附件的“评论”中了。我希望你不会觉得它质疑太过。我已经习惯在朋友中直接坦率地表达观点，这有时候会超过中国式礼貌可以容纳的程度。我这样做是因为我感觉你稿子中的内容十分有趣而且牵涉到的范围很广。我希望你能把评论看作是友好的。你也不用迟疑，请照直回复我。\end{bookquote}

% 原书第469页

作为一个后学，Paul如何质疑我的观点，都是我的荣幸。没有想到他既花时间给我提意见，还会想到要照顾我的感情。我在推出第一本书前，一直处于闭门造车的状态。我从来没有参加过学术会议，也没有给期刊投过稿，完全没有跟我博士答辩委员以外的学者有过交流。尽管我写的东西通过了博士答辩，拿到了学位，但实际上我对自己的论说是否能成立，其中是否有漏洞，心里面都没有底。Paul是我联系的第一位学者。能得到他毫无保留的评论，对我深入理解和评估自己的论点有很大帮助。由于Paul和我的研究高度契合，他所指出的问题都是有的放矢的真知灼见，我根据他的建议完善了底稿。也是读完我的底稿后，Paul得知我的研究还涉及了康熙朝的储位之争，他专门给我提到《北京纪事》第3卷中有一些相关内容，同时他还给发来了第4卷（目前尚未出版）的部分翻译底稿。他告诉我那里面也有一些相关内容。Paul知道储位之争是汉语语言圈学者特别关注的一个问题（英文学术圈对此问题不是很关注），所以他让我看看。我记得我当时快速翻看过一遍，并没有深入。我在读博士期间对储位之争问题还算有过仔细的了解，对各派的论说也大致明白。我当时的结论是这个热门话题已经被前人研究透了，没有继续研究的可能了。在这个固有想法的驱使下，外加我当时还在忙着出书，就没有真正静下心来细读Paul发给我的文档。


% 原书第470–471页

一直到2024年夏，《北京纪事》第3卷出版，我都没有想过要仔细看看里面的内容。我当时只是想参加Paul定在12月的网上活动，想争取在网上跟他打个招呼，并在活动中公开感谢他在这个翻译项目中的贡献。但出乎意料的是，11月底，中西研究e-mail群中一位老师发邮件告知大家Paul去世的消息。也是在这时，我才从纪念文章中仔细了解了一些Paul的生平。Paul是澳大利亚人，1937年1月19日出生于墨尔本，一生也主要在澳大利亚生活和做研究。他的中文名是鲁保禄，1955年他加入了耶稣会，但后来没有选择做一位传教士，在1964年退出教会后，他转向了学术研究。也是得益于早年在耶稣会学校中的学习，拉丁文以及欧洲的各种语言文字是他的强项。他最先研究的是欧洲中世纪史，后来在博士阶段开始学习中文，转向了中国的天主教史研究。因为早年与耶稣会的缘分，Paul的研究中心始终围绕着在华的耶稣会传教士展开。《北京纪事》的英文翻译2002年立项，由澳门利氏学社资助。Paul是项目负责人。我猜想Paul最初尽管知道这个项目宏大，但大概也没有预计到这个项目会持续超过二十年。由于Paul给我发过原始的Word底稿，我还能从他们内部团队写下的行边批注，看到许多尚待解决的问题。《北京纪事》的内容涉及面极广，里面有些用字母拼写出来的词汇可能是某个欧洲方言，也有可能是中国某个地方的方言，还有可能是满语。每一个小细节，都要耗费大量的时间和精力去破译。根据与Paul一起负责《北京纪事》的德国学者Claudia von Collani回忆（文章见Monumenta Serica 2025 年第 1 期），Paul从 2000 年开始就在着手撰写全面梳理中国礼仪之争历史的专著。可惜由于《北京纪事》占用了他大量的时间和精力，在其2024年11月20日去世前，预计三卷本的中国礼仪之争，他仅完成了第一卷，第二卷只写了一部分。我在感叹世事无常之余，又翻了一遍Paul给我发的众多电邮。这时我又看到了他给我说起关于储位之争的电邮。在这样的情况下，我开始细读《北京纪事》第3卷。很快我就发现里面确实如Paul所言，有新东西。


% 原书第471页

我最初的想法是写一篇论文，表示对Paul的纪念。Paul晚年花费大量心血翻译和注解《北京纪事》。对他而言，最大的欣慰就是有学者用到这份原始材料，能在此之上推出新的研究成果。因此我琢磨着基于《北京纪事》第3卷，写出一篇关于储位之争的论文。我想这样是对Paul最好的纪念。但很快我就发现这个问题由于前人的研究众多，从任何一个线头入手梳理，都有很多相关问题要回答。尽管我试着在很多地方进行概括性的书写，但还是很快就写了两万多字。这时论文还没有写完。


% 原书第471–472页

在写这篇论文前不久我正好投了一篇英文论文。投稿时我就感叹英文世界的大趋势是人文方面的期刊论文都在缩减字数。好多我熟悉的期刊都已经缩减到正文不得超过八千字、注释不得超过两千字的规模。我好不容易发现一个心仪的期刊的字数上限是一万两千字，但等我把稿件各方面修改好，准备正式投稿之时，发现该刊竟然已经改了规则，把字数上限降到了一万字。由于之前看到过一些比较长的中文论文，因而我对中文论文字数限制还抱有期待。可我在网上查了一圈我知道的中文期刊，发现都没有明确标出字数限制。我于是发微信问了在国内大学任教的朋友。他们告诉我国内期刊的字数按规定也就是在两万字左右，同时一位朋友坦率地告诉我，我看到的长文章是特例，不适合我这个“年轻人”。


% 原书第472页

当时我认为自己真看到了一些新材料，切实感觉到了敝帚自珍的含义。如果把论文缩减到两万字左右，很舍不得，心有不甘。而且之前为了照顾论文长度，还有很多内容根本未能放开手脚去书写。权衡再三，我想我只有写成专著（monograph），才能畅所欲言，没有限制。佛教中说起心动念都是缘。我想研究、写作说不定也应该追寻那心动一刻的缘分。为了避免自己遇到困难而变心，半途而废，也避免被其他邀约分心，我赶紧写了一份提纲，在2025年春节过后，提前去找出版社询问意向，希望赶紧把事情定下来。好在中华书局的编辑在看了目录、提纲和我已经写好的样章（现在第二章中的主要内容）后，马上就立了项，这件事情也就算定了下来。


% 原书第473页

2025年3月底，我回国参加我上一本书的一些活动时，有幸见到了本书的责任编辑孟庆媛。我们约在成都文殊院旁的茶馆见面。当天的成都阳光和煦，茶馆自然也是人声鼎沸，热闹非凡。传统喝盖碗茶的地方已经没有了位置，我们只得在抄写经书的院子中找到了一张桌子坐下，但依然需要提高嗓门才能交流。当天我给孟编辑讲了我看到的材料，以及打算怎么写，会有哪些写作的可能等等。孟编辑听得很有兴致，我感觉讲到关键处，她的眼中有光。西方写作中流传着一种说法，脑中的想法要自己能讲出来让别人明白，才算自己搞明白了。我给孟编辑讲了几个小时，感觉她都听得兴致盎然，一直能跟随我的思路，不断地问我问题，这让我感觉到我的叙事思路是有吸引力的。这大大增加了我完成整个题目的信心。


% 原书第473–474页

在我有了最初的想法后，就写了前言给我的大学同学洪晓敏和史翔看过，同时跟他们聊了我的想法。两位不在学界之中，又都是饱读诗书之人，他们可以从非学术视角给我提供意见。由于是老同学，我知道我的想法无论他们觉得好还是不好，都会直言相告。后来在确定书名时，我也最先去问他们的第一感觉。我到美国不久就认识的老友唐国靖也跟我聊了许多这本书中的问题。国靖在美国拿的是数学博士，但一直忧国忧民，对历史尤其是基督教在世界的传播史感兴趣。我们经常一起结伴外出滑雪，在开长途车时常常通过争论各种历史社会问题来打发时间。他反复告诉我，写书不要只想着学界中的那些问题，要有社会意识。他说的其实是要有市场意识。好的人文作品也可以是一件好的产品，能够进入到大众的生活中。


% 原书第474页

动笔之后，我顺着新材料探索过去，发现线头越理越多。最初预计写四章，后来增加到了七章。在校改过程中，又陆陆续续扩充了许多内容。从最初预计的两百五十页，变成了四百页，又到了现在近五百页的规模。最后我坚决打住，把这本书的主体内容限制在了康熙朝。在此要感谢孟编辑对我多次改动的迁就。全书写完后，还有幸请到中国社会科学院的定宜庄老师，中国人民大学的胡祥雨老师、董建中老师，北京大学的毛亦可老师，复旦大学的邹怡老师，以及《文史》的执行编委鲁明老师阅读。他们从不同方面给我指出了错漏，提出了进一步完善的建议。对此我十分感谢。我想他们在拿到成书的时候，会看到他们提的意见确实帮助到了这本书。我预计以后读者还会发现书中的问题，也请不吝在互联网中指出，让这本书未来能像其他创新产品一样，有机会在更新时不断完善。


% 原书第474–475页

互联网时代的作者和以前很大的一个不同就是现在的作者能看到大量读者的反馈。以前一些身边的朋友问我，你老写那些有啥用啊？我没有答案。而最近聚会时，我对他们说，写书有点像发了一个很大的朋友圈，能持续不断地看到各种有意思的评论。本书写到最后时，我不知道如何结尾，反复写了数个版本都不满意。10月中有一天下午去接小孩放学，在等的时候，我拿着手机翻看《康熙的红票》的豆瓣短评，看到有读者在评论的最后留下了一句：“就算换了时空变了容颜，我依然记得你眼里的依恋。”我突然感觉似乎有所领悟。回到家后，我在网上搜索后发现原来这是一句歌词，于是还找到了这首歌的原唱万芳的版本来听。歌声中唱出的对时空转化的无奈之感帮我为全书的结束找到了方向。当晚我很快就写出了我自己满意的结尾。“两句三年得，一吟双泪流。”有了辗转反侧写不出的经历，也才有这写出来后的感怀。写书是一段人生经历，其中也有这样意想不到的瞬间。


% 原书第475页

全书定稿后，我在收拾《北京纪事》时，又看到了纪理安神父在书前写下的，每次都很打动我的“致读者”。他在执笔时已经知道他们北京耶稣会顶撞了教皇，预感到他写的记录将有可能永不见天日。在这样的情况下，他解释了为什么还要坚持写完：


% 原书第475页

\begin{bookquote}Nevertheless, I write for eternity, and since I write the truth it will not be lacking in usefulness for the wise reader, and this will satisfy his debt to the writer rather than praise which the latter has been taught to avoid.\end{bookquote}

% 原书第475–476页

\begin{bookquote}无论如何，我是写给永恒时光的。既然我写的都是真相，那么总有明智的读者会发现它的用处。这就够了。这比赞美本书作者更能让他为其付出而感到满足。赞美从来都不是他追求的。\end{bookquote}

% 原书第476页

纪理安神父所谓“明智的读者”，在我看来就是中国文化中的“知音”。贾岛“两句三年得”写得那么苦，也是写给“知音”的，还说“知音若不赏”，欣赏不了的话，他就“归卧故山秋”，不写了。期待知音大概是中西方文字写作者共有的期待。而纪理安的文字让我特别感动之处是其静待知音的信念，以及信念中又蕴含着一种飞蛾扑火般的执迷不悔。我每次看完他这段话，心中都在默默地说，我发现了用处，应该算是一位明智的读者。希望本书能成为我的证明。


% 原书第476页

最后还要感谢读者，坚持读到这里。书写出来，有人愿意看，是写作者最大的光荣。如果您还愿意在互联网中写下感想，那是我的荣幸，我期待拜读。


% 原书第476页

\begin{flushright}
孙立天\\[.5em]
2025年12月\\[.5em]
于纽约郊外家中
\end{flushright}
